\section{The equal-leg endpoint and its neighbors}\label{sec:iso}

\begin{theorem}
At $u=\sqrt2-1$, the classical solution has no positive second brake and
reaches a binary or triple collision in finite time.
\end{theorem}
\begin{proof}
Here $m_1=m_2=m=1/\sqrt2$ and reflection symmetry persists.  Write the base
bodies as $(-x,y_b),(x,y_b)$ and the apex as $(0,y_a)$, with
$h=y_a-y_b$.  Then
\[
\ddot x=-\frac{m}{4x^2}-\frac{x}{(x^2+h^2)^{3/2}}<0.
\]
Starting with $x_0=1/2$ and $\dot x(0)=0$ gives $\dot x<0$, excluding a
second brake.  While $x>0$, one has
$\ddot x\leq-m/(4x_0^2)$, whence
\[
x(t)\leq x_0-\frac{m}{8x_0^2}t^2.
\]
Thus collision occurs no later than $\sqrt{8x_0^3/m}=2^{1/4}$.
\end{proof}

To study nearby rational parameters, we use the Levi--Civita continuation
of the symmetric orbit as a comparison solution. A collision of a nearby
physical trajectory still terminates its classical motion.

Put $v=u/(\sqrt2-1)$, $(m,n)=(A,B)$, $g=q_2-q_1=w^2$, and
$G=q_3-C_{12}$. With $dt=|w|^2d\sigma$, write $z=w_\sigma$, $P=\dot G$,
and $e=|\dot g|^2/2-(m+n)/|g|$. The selected pair collision is $w=0$;
this is regular in the comparison equations, but it still terminates the
classical trajectory at $v=1$.

For $d_{31}=G+ng/(m+n)$, $d_{23}=G-mg/(m+n)$, and
$F=d_{23}/|d_{23}|^3-d_{31}/|d_{31}|^3$, the regular equations are
\[
\begin{aligned}
 w_\sigma&=z,&
 z_\sigma&={e\over2}w+{|w|^2\bar w\over2}F,&
 e_\sigma&=2\operatorname{Re}(wz\bar F),\\
 G_\sigma&=|w|^2P,&
 P_\sigma&=-|w|^2{m+n+1\over m+n}
 \left({m d_{31}\over|d_{31}|^3}+{n d_{23}\over|d_{23}|^3}\right).
\end{aligned}
\]
They preserve $2|z|^2-(m+n)-e|w|^2=0$ and are analytic at $w=0$
when the other two separations are positive. At $v=1$ the launch is
$(w,z,e,G,P)=(1,0,-\sqrt2,i/2,0)$; its only nonzero parameter
derivatives are
$m_v=-1/2$, $n_v=1/2$, and $G_{x,v}=1/(2\sqrt2)$.

\subsection{The initial arc and collision unfolding}

The first syzygy, or collinear configuration, is reached without a brake
for all tied parameters sufficiently near $v=1$. The relevant variational
argument also establishes the torque inequality used in
Section~\ref{sec:torque}. Write
$F_0(v,t)=\eta(v,t)-h(v,t)$ for the history-threshold difference defined
there, and let $\tau(v)$ be the first syzygy time. On the symmetric orbit,
$F_0(1,t)=0$. Differentiation of the Newtonian equations and the exact
tied launch gives an analytic function $g(t)=\partial_vF_0(1,t)$ with
\[
 g(0)={32+8\sqrt2\over7}>0.
\]
A validated orbit-and-tangent cover proves $g>6$ up to $t=10^{-3}$,
$g>4$ on $[10^{-3},0.43]$, and $g_t<-6$ from $0.43$ through the
unique transverse syzygy. The collinear identity for the torque threshold
gives $F_0(v,\tau(v))=0$, hence $g(\tau(1))=0$. Analytic division yields
\[
 F_0(v,t)=(v-1)(t-\tau(v))F_2(v,t),\qquad
 F_2(1,\tau(1))=g_t(\tau(1))<0.
\]
This factorization controls the terminal part of the arc, while compact
continuity controls its interior. Thus $F_0<0$ for $v<1$ sufficiently
close to one and $0<t<\tau(v)$. The same tangent cover proves
$r_{12}>r_{23}>r_{31}$ there. A pair angular momentum then remains
nonzero through syzygy, excluding a brake. The launch division and
orbit-and-tangent equations are checked by exact symbolic regression;
the full interval inputs are indexed in Appendix~\ref{sec:reproduction}.

At the first endpoint collision, an oriented $C^1$ Poincar\'e calculation
gives $z_r<-4/5$ and $\partial_v w_i<-30$. Thus the map
$(\sigma,v)\mapsto(w_r,w_i)$ has nonsingular derivative there. The nearby
members miss that collision, with
\[
 \min r_{12}=\chi^2(1-v)^2+o((1-v)^2),\qquad \chi^2>900.
\]
At the section $w_r=0$,
$\ell_{12}=-2w_i z_r>0$ for $v<1$ sufficiently close to one, whereas
its launch sign is negative. Thus the launch sign has already reversed
by this section; it cannot be assumed to persist through later encounters.

\subsection{Continuation to a terminal escape section}

Four comparison collisions must be controlled before the escape criterion
applies. Their regularized transversality and the nonzero brake residual
give a single punctured neighborhood on which the entire argument holds.

\begin{theorem}[Punctured near-isosceles nonperiodicity]
There exists $v_*<1$ such that every real tied member with
$v_*<v<1$ is nonperiodic.  It is collision-free and has no labelled brake
through a fixed LC section; thereafter it either ends in a classical binary
collision or has body 3 escape.  Consequently the same conclusion holds for
infinitely many primitive Pythagorean triples.
\end{theorem}

\begin{proof}
The exact LC energy constraint is
\[
 2|z|^2-(m+n)-e|w|^2=0.
\]
On the symmetric endpoint orbit it gives $2|z|^2=\sqrt2$ whenever $w=0$,
so every selected collision is transverse.  Pinned oriented $C^1$ Poincare
maps enclose four such zeros before $\sigma=7$ and put the fifth strictly
after $7$.  At each of the first four zeros, the tied normal derivative
$(w_i)_v$ is bounded away from zero.  Thus every Jacobian
\[
 {\partial(w_r,w_i)\over\partial(\sigma,v)}
\]
is nonsingular.  The inverse-function theorem isolates all four zeros, while
a 7000-step $C^0$ cover proves
\[
 r_{31}^2,r_{23}^2>{1\over500}
 \qquad(0\leq\sigma\leq7).
\]
Compactness outside the four collision boxes now gives a punctured
neighborhood in which every nearby tied orbit is collision-free through
$\sigma=7$.

At an ordinary point of this LC chart,
\[
 \dot g={2z\over\bar w},\qquad \dot G=P.
\]
Hence, after center-of-mass reduction, a labelled brake is equivalent to
$z=P=0$.  The polynomial residual
$\mathcal R_{\rm LC}=|z|^2+|P|^2$ remains regular at comparison collisions.
At the $\sigma=1/2$ interface the cover gives $G_y>1/10$, so the unique first
syzygy lies later.  Every complete interval tube from that overlapping
pre-syzygy interface to
$\sigma=7$ satisfies $\mathcal R_{\rm LC}>2$.  Compact continuity transfers
this inequality to the nearby classical orbits; the initial-arc result described above covers the initial segment.

At $\sigma=7$, take $\eta=4$ in Lemma~\ref{lem:escape}. The
same interval cover proves, with $M=m+n=\sqrt2$, $R=M/4$ and
$s=|G|-R$,
\[
 s>1,\qquad \dot\rho>2,\qquad
 \delta={\dot\rho^2\over2}-{M+1\over s}>{1\over20},
\]
and
\[
 -4-e-{\sqrt{2MR}\over\sqrt{2\delta}\,s^2}>2.
\]
All inequalities are strict and therefore persist in a tied neighborhood.
The lemma says that each future trajectory either suffers an
inner classical collision or escapes with uniformly positive outward radial
speed.  Neither possibility permits a later brake.  Rational Euclid
parameters are dense in the resulting open $u$ interval, giving infinitely
many primitive triples.
\end{proof}

The cutoff $v_*$ is existential. No numerical width is claimed for this
endpoint neighborhood. An effective bound would identify a corresponding
explicit range of integer triangles.
